\documentclass[11pt]{article}
\usepackage[a4paper,margin=25mm]{geometry}
\usepackage[T1]{fontenc}
\usepackage{lmodern,booktabs,multirow,xcolor}
\definecolor{navy}{HTML}{172B43}
\pagestyle{empty}
\begin{document}
\noindent{\small\sffamily\color{navy} DDA5001 / SYNTHETIC REGRESSION STUDY}\par
\vspace{8mm}
\noindent{\LARGE\bfseries Baseline and linear regression}\par
\vspace{3mm}
\noindent A concise comparison on two synthetic datasets.\par
\vspace{10mm}
\begin{center}
\renewcommand{\arraystretch}{1.25}
\setlength{\tabcolsep}{8pt}
\begin{tabular}{llrrrr}
\toprule
\multirow{2}{*}{Dataset} & \multirow{2}{*}{Model} &
\multicolumn{2}{c}{Training} & \multicolumn{2}{c}{Held-out test} \\
\cmidrule(lr){3-4}\cmidrule(lr){5-6}
 & & MSE & MAE & MSE & MAE \\
\midrule
\multirow{2}{*}{Low noise} & Mean & 1.4221 & 1.0051 & 1.5476 & 1.1372 \\
 & Linear & 0.0154 & 0.1012 & 0.0133 & 0.1011 \\
\addlinespace[6pt]
\multirow{2}{*}{Higher noise} & Mean & 1.6185 & 1.0422 & 1.8640 & 1.2140 \\
 & Linear & 0.1706 & 0.3373 & 0.1474 & 0.3371 \\
\bottomrule
\end{tabular}
\end{center}
\vspace{4mm}
\noindent\textbf{Table 1.} Single-split results on synthetic data
($y=2x+1+\epsilon$). Gaussian noise standard deviations are 0.15 (low)
and 0.50 (higher). Each dataset uses 30 training and 10 test observations,
with seed 42. MSE is mean squared error; MAE is mean absolute error.
Lower values indicate smaller prediction errors.\par
\vspace{7mm}
\noindent\textbf{Models.} Mean predicts the training-target mean. Linear fits
a slope and intercept using training data only.\par
\vspace{5mm}
\noindent\textbf{Interpretation.} The fitted line has lower held-out errors
than the constant baseline on both datasets. These small synthetic examples
illustrate a reporting workflow; they do not establish general performance.\par
\vspace{5mm}
\noindent\textbf{Source.} Values copied from \texttt{results.csv} at four
decimal places. No uncertainty estimates are implied.\par
\end{document}
